\section{Extremal lattice sections and design moments}
\label{sec:sections}

Dimension 43 counts minimal vectors perpendicular to a five-dimensional
section; dimension 45 projects three fibres of a three-dimensional
section. Complete projection domains and design moments determine the
required statistics without enumerating the ambient shell.

\subsection{Mother shells and moment equations}
An extremal even unimodular lattice of rank 48 has minimum squared norm
six and theta series
\[
E_4^6-1440E_4^3\Delta+125280\Delta^2
=1+52416000q^3+O(q^4).
\]
Here $E_4=1+240q+2160q^2+\cdots$ and
$\Delta=q-24q^2+252q^3+\cdots$; the constant term and vanishing first
two coefficients determine this weight-24 form~\cite{conway1999sphere}.
Its minimal shell $X$ is an 11-design~\cite[Theorem 4.1]{nebe-designs}.

For section vectors $t_i$ with Gram $H$, let $y_i=x\cdot t_i$ and
$f(y)=|\{x\in X:y(x)=y\}|$. Integrality, $y^{\mathsf T}H^{-1}y\le6$
and products with short section vectors restrict the signature domain.
For $2m\le10$, the moments are
\begin{equation}\label{eq:section-moment}
\sum_{x\in X}\prod_{\ell=1}^{2m}(x\cdot t_{i_\ell})
=\frac{|X|6^m}{48\cdot50\cdots(48+2m-2)}
\sum_{\mathcal P}\prod_{(a,b)\in\mathcal P}H_{i_ai_b},
\end{equation}
where $\mathcal P$ pairs labelled positions; degree zero gives $|X|$.
These include even total degrees with odd individual exponents.
For a moment system $Av=b$, a rational identity $A^{\mathsf T}w=c$
proves $c^{\mathsf T}v=w^{\mathsf T}b$. Averaging must preserve the
domain, moments and target, but need not extend to an ambient
lattice automorphism.

\subsection{Dimension 43: the fixed child and eight coordinate cuts}
In $P_{48n}$~\cite{latticecatalogue}, supplied integer rows $T$ embed
$\sqrt3E_8$: the ambient Gram $G$ satisfies
$(TGT^{\mathsf T})_{ij}=3\alpha_i\cdot\alpha_j$, where
\[
\alpha_1=\tfrac12(1,-1,-1,-1,-1,-1,-1,1),\quad
\alpha_2=e_1+e_2,\quad
\alpha_j=e_{j-1}-e_{j-2}\quad(3\le j\le8).
\]
The first five rows specify the $\sqrt3D_5$ child $U$.
This embedded-parent approach belongs to the cross-section framework
of Dorofeev--Sun--Wang~\cite{dorofeev}.

\begin{proposition}[$E_8$-extension count]\label{prop:section43}
Let $L$ be an extremal even unimodular lattice of rank 48 with an
isometric embedding $\iota:\sqrt3E_8\hookrightarrow L$, and let $X_L$
be its squared-norm-six shell. The span $U_L$ of the first five embedded
simple roots, and every $W(E_8)$-conjugate of this child, satisfies
\[
|X_L\cap U_L^\perp|=2553792.
\]
No ambient Weyl action or primitivity assumption is required.
The supplied embedding gives $K(43)\ge2553792$.
\end{proposition}

The parent projection is $\lambda/\sqrt3$ with $\lambda\in E_8$,
by integrality and self-duality. Apart from the 240 singleton
exceptions $\lambda=3r$, the minimum gives
$\|\lambda\|^2\le18$ and $|\lambda\cdot r|\le3$ for every root $r$.
Exact enumeration gives $1,240,2160,6720,17280$ signatures at squared
norms $0,2,4,6,8$: 26641 signatures including exceptions.
The radial moments force Weyl-averaged multiplicities
\[
(a_0,a_2,a_4,a_6,a_8)=(1092000,116235,9180,495,63/4).
\]
The child's perpendicular space meets these orbits in $1,12,6,24,0$
signatures, with twelve exceptions, giving the average count
$1092000+12(116235)+6(9180)+24(495)+12=2553792$.

For the fixed child, the group $\langle W(D_5),W(A_3),-I\rangle$
preserves the target and has 43 domain orbits; $A_3$ consists of the
roots perpendicular to the child. The certificate gives
$A^{\mathsf T}w=c$ and $w^{\mathsf T}b=2553792$ for regenerated
moment and fixed-exception rows. This proves the proposition without
assuming ambient symmetry. The complete domain is essential: the
orthogonality condition belongs in $c$, not in prior column deletion.
Normalizing by $\sqrt6$ gives a kissing configuration, since distinct
minimal vectors have product at most three. The comparison has
2545056 points from an isometric $D_5$: the gain is in its embedding.

Extend $\iota$ linearly and put
$U_k=\iota(\sqrt3\operatorname{span}\{e_1,\ldots,e_k\})$,
$1\le k\le8$. With $m=8-k$, the eight fixed counts are
\begin{align*}
|X_L\cap U_k^\perp|={}&1092000+18360m+464944\binom m2
+11880\binom m3\\
&+146880\binom m4+2520\binom m5+31680\binom m6.
\end{align*}
These concern the specified cuts and their Weyl conjugates, not
arbitrary sections. $U_1$ is an axis, not a $D_1$ root system;
the other cuts are coordinate $D_k$ spans.

For completeness, in doubled coordinates $u=2\lambda$, group the
nonexceptional domain by permutations within the first $k$ and last
$8-k$ positions and even sign changes. The polynomials of degree at most ten
$m_a(u_1^2,\ldots,u_k^2)m_b(u_{k+1}^2,\ldots,u_8^2)$,
$|a|+|b|\le5$, use monomial symmetric functions $m_a$ and moments
of Gram $12I_8$. The eight orbit matrices have full column ranks
$25,34,39,41,39,34,25,13$, certified by nonzero integer minors.
Thus actual orbit masses equal the Weyl-averaged masses $|O|a_r$.
Counting target signatures gives the formula, whereas radial moments
alone would give only an average over cuts.
The Weyl map
$\tfrac12P_{5,8}\operatorname{diag}(J_4-2I_4,J_4-2I_4)$
sends the specified $D_5$ span to the first five coordinate axes.
Here $P_{5,8}$ exchanges coordinates five and eight.

\begin{center}
\begin{tabular}{r|rrrr}
$k$&1&2&3&4\\\hline
$|X_L\cap U_k^\perp|$&16815624&10663920&6688960&4149504\\[2pt]
$k$&5&6&7&8\\\hline
$|X_L\cap U_k^\perp|$&2553792&1593664&1110360&1092000
\end{tabular}
\end{center}

\subsection{Dimension 45: an exact fibre identity}
Appendix~\ref{app:ambient} proves minimum six for the specified
quadratic-residue even neighbour $\Lambda$. More generally, the
following argument applies to any extremal even unimodular rank-48
lattice with a section of Gram
$H=8I_3-2J_3$, $H^{-1}=(I_3+J_3)/8$, without primitivity.

The section vectors $t_1,t_2,t_3,-t_1-t_2-t_3$ and their negatives
give eight singleton signatures
$E=\{\pm He_1,\pm He_2,\pm He_3,\pm H(1,1,1)^{\mathsf T}\}$.
Every other signature lies in
\begin{equation}\label{eq:domain45}
D_0=\{y\in\mathbb Z^3:|y_i|\le3,\ |y_1+y_2+y_3|\le3\}.
\end{equation}
Its 231 signatures have projection norm at most $9/2$.
Subtracting $E$ and pairing antipodes gives 161 moment rows and
116 variables, with $A_{a,[y]}=(2-\mathbf1_{y=0})y^a$ and
$b_a=M_a-\sum_{y\in E}y^a$.

\begin{proposition}[Exact fibre identity]\label{prop:moment45}
Under these hypotheses,
\[
f(e_1)+f(e_2)+f(e_3)=7377408+9f(-2,-2,3).
\]
\end{proposition}
\begin{proof}
Give $c$ coefficients $1$ on $[e_i]$, $-9$ on $[(-2,-2,3)]$ and
zero elsewhere. The rational multiplier in \texttt{dual.json}
(Appendix~\ref{app:data}) satisfies
$c-A^{\mathsf T}w\ge0$, $w^{\mathsf T}b=7377408$.
Since $t_i+t_j$ has squared norm eight, the minimum also gives
$|y_i+y_j|\le4$. With $y_4=-\sum_{i=1}^3y_i$, this removes the 30
permutations of $(3,3,-3,-3)$ and $(3,2,-3,-2)$.
All seven positive slack entries lie on these excluded orbits.
Thus $A_D^{\mathsf T}w=c_D$ on the remaining 101 antipodal orbits,
sharpening the original lower certificate to equality.
\end{proof}

Projecting the three fibres leaves squared norm $23/4$ and products
at most $11/4$ within one fibre and $23/8$ between fibres. This
excludes collisions and gives a kissing configuration after
normalization. Thus $m$ distinct witnesses in the last fibre give
$K(45)\ge7377408+9m$; its complete size gives an exact cardinality.

\subsection{The anchor graph and the certified count}
For a norm-eight anchor $p$, put $A_p=\{x\in X:x\cdot p=4\}$.
The minimum gives $|x\cdot p|\le4$. The filter
$P(t)=t^2(t^2-1)(t^2-4)(t^2-9)$ has $P(4)=20160$ and shell sum
90961920. Antipodality therefore gives $|A_p|=2256$.
Pair $x$ with $p-x$. Distinct pairs have endpoint products in
$\{1,2,3\}$; exchanging an endpoint replaces the product by four
minus itself. Products one or three thus define a graph on 1128 pairs.

Filtering the mixed design moments through degree ten gives
$\sum_{x\in A_p}(z\cdot x)^2=192\|z\|^2$ for $z\perp p$.
For $z=r-p/2$, its own pair contributes 32 and each adjacent pair
contributes two. Hence $768=32+2\deg(\bar r)$: every such graph is
368-regular.

For nonadjacent vertices represented by $r,s$, the section
$(-s,s-p,p-r)$ has Gram $H$. Each vertex of
$N(\bar r)\setminus N(\bar s)$ has a unique endpoint $a$ with
$r\cdot a=3$, and $w=p-a$ has signature $(-2,-2,3)$.
Conversely, that signature forces $p\cdot w=4$, $r\cdot w=1$ and
$s\cdot w=2$, recovering exactly that endpoint. The complete fibre
therefore has size $368-c$, where $c=|N(\bar r)\cap N(\bar s)|$,
and the projected count is $7380720-9c$.
Every ordered section of Gram $H$ has this representation:
$p=-t_1-t_2$, $s=-t_1$, $r=-t_1-t_2-t_3$.

Choose one endpoint $r_i$ per pair and set $u_i=r_i-p/2$.
The signed matrix $S$ has zero diagonal and entries
$S_{ij}=u_i\cdot u_j\in\{0,\pm1\}$ for $i\ne j$. The filtered
second moment gives
\[
\sum_i u_i\otimes u_i=96I_{p^\perp},\qquad
S^2=88S+368I,\qquad
\operatorname{Spec}(S)=\{92^{(47)},(-4)^{(1081)}\}.
\]
Indeed, the Gram $4I+S$ has rank 47 and squares to 96 times itself.
Signed common-neighbour counts satisfy
$c_{ij}^{\pm}=(c_{ij}\pm88S_{ij})/2$; thus codegrees are even and
projected counts are divisible by 18.
The unsigned adjacency matrix satisfies $A+16I\succeq0$ as the Gram
of $u_i\otimes u_i$; this does not imply a two-valued unsigned spectrum.
Four signed $u_i$ cannot have all mutual products $-1$: their sum
would be a lattice vector of squared norm four.

The selected anchor has numerator shape $(3^6,1^{42})$ in coordinates
scaled by $1/\sqrt{12}$. Its 2256 certified endpoints are complete.
All 428076 nonadjacent pairs have codegrees 70 to 148, with a unique
pair attaining 70. Its fibre therefore has exactly 298 points,
improving on the earlier 292-witness Pless-lattice construction.

\begin{corollary}[The 45-dimensional construction]\label{cor:section45}
The specified construction has exactly
$7377408+9(298)=7380090$ points, hence $K(45)\ge7380090$.
This is the largest projected count for this fixed anchor and Gram
$H$, not an optimum over other anchors, lattices or section Grams.
\end{corollary}
